\documentclass[11pt,letterpaper]{article}
\usepackage[T1]{fontenc}
\usepackage{lmodern}
\usepackage[margin=0.95in]{geometry}
\usepackage{amsmath,amssymb,amsthm,mathtools,booktabs,longtable,array}
\usepackage{microtype,xcolor,hyperref,fancyhdr,listings}
\definecolor{inkblue}{RGB}{25,55,90}
\hypersetup{colorlinks=true,linkcolor=inkblue,urlcolor=inkblue,citecolor=inkblue,
 pdftitle={Twelve-variable Diophantine forms of ZFC consistency},pdfauthor={Research synthesis for Vladimir Reshetnikov}}
\pagestyle{fancy}\fancyhf{}\fancyhead[L]{\small Diophantine forms of ZFC consistency}
\fancyhead[R]{\small 22 October 2026}\fancyfoot[C]{\thepage}
\setlength{\headheight}{14pt}
\setlength{\emergencystretch}{2em}
\newtheorem{theorem}{Theorem}[section]
\newtheorem{lemma}[theorem]{Lemma}
\newtheorem{proposition}[theorem]{Proposition}
\theoremstyle{remark}\newtheorem{remark}[theorem]{Remark}
\newcommand{\Z}{\mathbb Z}\newcommand{\N}{\mathbb N}\newcommand{\pos}{\mathbb Z_{>0}}
\newcommand{\Con}{\operatorname{Con}}\newcommand{\Res}{\operatorname{Res}}
\newcommand{\pc}{\operatorname{pc}}\newcommand{\PS}{\mathcal P}
\newcommand{\pin}{\texttt{f0ae4e0991ca44a7bb4e5671d42e33d775e48035}}
\newcommand{\pathtext}[1]{\texttt{\detokenize{#1}}}
\newcommand{\IntegerDegree}{2993}
\newcommand{\IntegerOps}{169}
\lstset{basicstyle=\ttfamily\small,breaklines=true,columns=fullflexible,frame=single,
 rulecolor=\color{black!20},showstringspaces=false}
\title{\textbf{Twelve-variable Diophantine forms\\of ZFC consistency}\\[0.6em]
 \large Monic elimination, signed domains, and smaller fixed coefficients}
\author{Research synthesis prepared for Vladimir Reshetnikov\\\small ProveIt integration package}
\date{22 October 2026}
\begin{document}
\maketitle
\begin{abstract}
We combine the existing blank-input K215 compiler, weighted square-root elimination,
and signed-domain arguments with a new elimination based on algebraic integers.
The resulting explicit polynomial has twelve positive integer unknowns, degree
1529, and a division-free evaluator with 117 binary arithmetic operations. Its
three-equation presentation has maximum degree 1520 and costs 112 operations.
An intermediate thirteen-variable polynomial has degree 319, improving the
reviewed thirteen-variable degree 515. A signed version uses 26 unrestricted
integer unknowns, improving 29. A separate combination gives a 59-operation,
seven-equation system in sixteen positive unknowns. Finally, a strict four-cell
compiler uses 5,252 tiles and base $2^{5^{16}}$, reducing the binary length of
the base-minus-one coefficient by a factor of 25 from the reviewed T4 recipe.
All constructions have solutions exactly when K215 halts on its blank tape,
and hence their unsolvability is equivalent in ZFC to syntactic $\Con(\mathrm{ZFC})$.
The new mathematics is proved here and accompanied by finite, reproducible
checks. It is not a new Lean or Rocq formalization and makes no global minimality claim.
\end{abstract}

\section{Results, comparison, and the precise claim}

\subsection{The outcome}
The main new equation is
\begin{equation}\label{eq:main}
 \boxed{\PS_{12}=N_kN_t(1+\mathcal R_{12})-1=0,}
\end{equation}
in the twelve positive integer unknowns
\begin{equation}\label{eq:variables}
 J,F,\alpha,z,h,i,T,s,w,\eta,\zeta,u.
\end{equation}
Every symbol on the right of \eqref{eq:main} is an explicitly defined polynomial
abbreviation. Section~\ref{sec:equation} defines them, including a division-free
formula for $\mathcal R_{12}$. Four fixed integer coefficients are specified by
a terminating recipe in Section~\ref{sec:coefficients}. There is no unspecified
machine index, no unknown coefficient, and no variable exponent in the equation.

\begin{theorem}\label{thm:main}
For the fixed coefficient tuple of Section~\ref{sec:coefficients},
\[
 \mathrm{ZFC}\ \vdash\
 \bigl[\neg\exists\boldsymbol{x}\in\pos^{12}\,
       \PS_{12}(\boldsymbol{x})=0\bigr]
 \ \longleftrightarrow\ \Con(\mathrm{ZFC}).
\]
The corresponding three-equation system $N_k=N_t=1$, $\mathcal R_{12}=0$
has the same existence condition. The polynomial in 26 unrestricted integers
defined in Section~\ref{sec:integers} has the analogous equivalence over $\Z$.
\end{theorem}

The proof is a composition of ordinary mathematical arguments formalizable in
ZFC. In particular, it does not infer an internal ZFC theorem merely from the
soundness of a proof assistant. The inherited K215 calibration, arithmetic
compiler, and new algebraic reductions all admit the usual finite coding
within set theory. The machine endpoint has existing Lean and Rocq evidence;
the new Diophantine bridge does not acquire that formal status by composition.

\subsection{A comparison with the reviewed snapshot}
The review is pinned to \pin. This snapshot contains the seven wave-2 packages
numbered 09--15. Their reconciliation note is included as
\pathtext{reference/comparison_at_pin.md}; its statistics are the packages'
stated figures, not new independent rebuilds of all seven packages.
The table compares matching objectives rather than declaring one universal
notion of simplicity.

\begin{center}\small
\begin{tabular}{>{\raggedright\arraybackslash}p{3.4cm}rrrrp{2.2cm}}
\toprule
Construction & Unknowns & Eqs. & Degree & Ops. & Domain\\
\midrule
Reviewed 09 low weight &13&1&515&85&positive\\
\textbf{New main} &\textbf{12}&1&1529&117&positive\\
New main system &12&3&1520&112&positive\\
\textbf{New intermediate} &13&1&\textbf{319}&91&positive\\
New intermediate system &13&4&288&83&positive\\
Reviewed integer best (14) &29&1&188&121&unrestricted $\Z$\\
\textbf{New integer form} &\textbf{26}&1&\IntegerDegree&\IntegerOps&unrestricted $\Z$\\
Reviewed retained system (09) &16&7&24&60&positive\\
\textbf{New retained system} &16&7&24&\textbf{59}&positive\\
Reviewed minimum-ops polynomial &15&1&155&67&positive\\
\bottomrule
\end{tabular}
\end{center}

Here ``degree'' for a system means the maximum total degree of its residual
polynomials. An operation is a binary $+$, $-$, or multiplication, including
multiplication by a fixed numeral. Reusing an already computed register is
free; supplying a fixed coefficient is free; equality tests are not counted
as arithmetic gates. Polynomial abbreviations do not introduce unknowns.
All new circuit counts have been checked on the emitted live gate arrays.

The twelve-variable result improves the reviewed explicit variable count, at
the cost of higher degree and more arithmetic. The thirteen-variable result
improves degree at the same variable count. The 59-operation result improves
the system count. We do not find a polynomial below 67 operations. The existing
32-variable quadratic lift, degree-22 fifteen-variable systems, and other
low-degree alternatives remain useful for their respective objectives.

There is also a separate coefficient improvement:
\begin{center}
\begin{tabular}{lrrr}\toprule
Recipe & Tiles & Base $B_{\mathrm{base}}$ & Bits in $B_{\mathrm{base}}-1$\\\midrule
Reviewed T4 (11)&37,075&$2^{5^{18}}$&$3,814,697,265,625$\\
New strict four-cell&5,252&$2^{5^{16}}$&$152,587,890,625$\\\bottomrule
\end{tabular}
\end{center}
Even the new coefficients are enormous. They are specified exactly, not
expanded. The bit count for $B_{\mathrm{base}}-1$ is about 19.1 decimal GB;
the coefficient recipe is the practical representation.

\subsection{What was combined}
Package 09 contributes weighted elimination of the strong, first, and auxiliary
Pell roots. Package 12 supplies a signed main-root and auxiliary-quotient
interface, together with the transport-coordinate shift. Package 14 shows
how to retain only seven positive coordinates by choosing an odd multiplier
and a square power scale. Package 11 supplies the four-cell local language,
native compiler layout, and a much smaller coefficient tuple.

The new step removes the main congruence quotient together with the old
combined-root coordinate. A monic polynomial makes the replacement root an
algebraic integer. A short norm argument forces it to be rational and hence
integral. A separate Pell lemma proves that the remaining rational auxiliary
ordinate is automatically integral. These two facts are what make the
resultant elimination sound; merely taking a resultant would not suffice.

Classical universal constructions can use fewer unknowns: the repository's
formal development includes an existence theorem for nine unknowns, while
eleven-variable integer representations are discussed in modern work
\cite{Jones,Bayer}. Our claim is a smaller concrete specialization of the
reviewed native-compiler constructions with explicit coefficients, equations,
degree, and evaluator. It is not a new universal variable bound.

\section{The inherited arithmetic interface}\label{sec:interface}

Write the four fixed coefficients as
\[
 \beta=B_{\mathrm{base}}-1,\qquad K_0=D_C+B_{\mathrm{base}}D_R,
 \qquad M_C,\qquad M_F^*=M_F+B_{\mathrm{base}}-1.
\]
The star is essential: the polynomial uses the shifted mask, not $M_F$.
The compiler establishes
\begin{gather*}
 0<M_C,M_F<B_{\mathrm{base}}-1,\qquad
 M_C\equiv2\pmod4,\qquad M_F\equiv4\pmod8,\\
 \pc(M_C)+\pc(M_F)=\log_2 B_{\mathrm{base}}.
\end{gather*}
All four ports must be taken from one whole authenticated recipe.

For positive supplied coordinates, define
\begin{align}
 q&=\beta J+1,&C&=q-F-\alpha,\nonumber\\
 X&=wq,&Y&=sq^3,&E&=XY,\nonumber\\
 k&=\eta+\zeta,&c&=kY+\eta,&a&=E+Y,\label{eq:common1}\\
 H&=4a+3,&\Delta&=a^2+H=(a+2)^2-1,\nonumber\\
 G&=(q-1)(q-F)+\alpha,&
 R&=G(q^2-1)+(M_C+qM_F^*)J,\nonumber\\
 L_0&=E(kY),& A&=1+L_0(L_0+k),\nonumber\\
 N_k&=k-hE-R,&N_t&=(K_0+w)C+(q-F)-z(q-1).\label{eq:outer}
\end{align}
The letter $A$ here denotes the first Pell radicand. The usual Pell sequence
parameter is $a+2$, which we will sometimes call $A_0$.

\begin{samepage}
The normalized fifteen-coordinate interface restores four further witnesses
$f,\tau,y,\gamma$ and has equations
\begin{align}
 \tau^2&=A,\label{eq:oldfirst}\\
 (X+ac+H\gamma)^2-\Delta c^2&=1,\label{eq:oldmain}\\
 f^2-\Delta i^2c^4&=1,\label{eq:oldstrong}\\
 QV_f^2-(Q-1)y^2&=1,\label{eq:oldaux}\\
 N_k=N_t&=1,\label{eq:oldouter}
\end{align}
\end{samepage}
where $Q=(\Delta i c^2)^2$ and
$V_f=c(Tf-1)-Rf^2$. There are fifteen supplied coordinates before the
eliminations: the twelve in \eqref{eq:variables}, with $u$ removed and
$f,\tau,y,\gamma$ inserted.

The inherited compiler proves that this system has a positive solution
exactly when blank K215 halts. Its signed extension permits the appropriate
restored roots, $T$, and $\gamma$ to be signed. One way to apply the ordinary
strong-norm version to \eqref{eq:oldstrong} is to put
\begin{equation}\label{eq:ordinaryconvert}
 i_{\rm ordinary}=\Delta i.
\end{equation}
Then $(i_{\rm ordinary}c^2)^2=\Delta(f^2-1)$, which is exactly its strong
equation. The conversion is a proof step; it introduces no new supplied variable.

For reference, the main logical steps in the inherited argument are as follows.
Weak transport gives $C\geq0$ even when $N_k=N_t=-1$. The authentic masks yield
\begin{equation}\label{eq:bootstrap}
 q\geq16,\quad q(q^2-1)<R<q^4-2,
 \quad E\geq q^4>R+2,\quad Y\geq4096.
\end{equation}
The first and main Pell equations classify $k=2\psi_{2XY^2+1}(n)$ and
$c=\psi_{a+2}(p)$. The strong equation and signed auxiliary congruences recover
$p=R$. The no-wrap estimate and a Pell ratio exclude the negative outer
unit signs. The main-root congruence then gives $X=2^R$, with the negative
main root excluded by the signed argument. Half-binomial extraction restores
the population condition and the native masks. Carry-free native convolution
recovers a periodic local computation tableau containing Start. A finite
strip around that Start gives a genuine blank halting run.

Conversely, any finite halting run admits independent spatial and temporal
padding. Free dummy bits tune its actual packed rank, and the positive Pell
converse supplies all remaining coordinates. Detailed inherited proofs are
included under \pathtext{reference/}; the new compiler adapter is proved in
Section~\ref{sec:coefficients}.

\section{The twelve-variable equation in full}\label{sec:equation}

In addition to \eqref{eq:common1}--\eqref{eq:outer}, put
\begin{align}
 I&=ic^2,&D&=1+\Delta I^2,&Q&=\Delta(D-1),&d&=Q-1,\nonumber\\
 v_0&=cT,&b_0&=c+RD,\nonumber\\
 P&=QD v_0^2,&Z&=Qb_0^2,&e&=P+Z-1,&a_0&=PZ,\label{eq:weighted}\\
 M&=2(kI^2+T^2),&&N=dM^2A.\nonumber
\end{align}
In particular $Q=(\Delta i c^2)^2>1$. Define a monic polynomial in a formal
indeterminate $v$:
\begin{align}
 \mathcal H_8(v)={}&
 \bigl[v^4-2d(N+e)v^2+d^2((N-e)^2+4a_0)\bigr]^2\nonumber\\
 &\quad-16a_0d^2\bigl[v^2+d(N-e)\bigr]^2.\label{eq:h8}
\end{align}
The indeterminate $v$ is eliminated; it is not a thirteenth unknown. Put
\begin{equation}\label{eq:mainquadratic}
 B_*=1+\Delta c^2,\qquad W=X+ac+Hu,\qquad p=Hc,
\end{equation}
and let
\begin{equation}\label{eq:resultant}
 \mathcal R_{12}=\Res_v\bigl(\mathcal H_8(v),(W-pv)^2-B_*\bigr).
\end{equation}
Equation \eqref{eq:main}, with these definitions, is the promised polynomial.

\subsection{A division-free evaluator}
The notation ``resultant'' in \eqref{eq:resultant} is only a concise
specification. The following straight-line formulas evaluate exactly the same
integer polynomial using only addition, subtraction, and multiplication:
\begin{align}
 r_0&=W^2+B_*,& r_1&=2W,& t_0&=dp^2,&t_1&=t_0^2,\nonumber\\
 S&=N-e,& \ell&=2t_0(N+e),&v_c&=t_1(S^2+4a_0),\nonumber\\
 v_R&=r_0^2+B_*r_1^2-\ell r_0+v_c,
 &v_I&=r_1(2r_0-\ell),\nonumber\\
 u_R&=r_0+t_0S,&\lambda&=16a_0t_1,\label{eq:normcircuit}\\
 h_R&=v_R^2+B_*v_I^2-\lambda(u_R^2+B_*r_1^2),\nonumber\\
 h_I&=2(v_Rv_I-\lambda u_Rr_1),\nonumber\\
 \mathcal R_{12}&=h_R^2-B_*h_I^2.\nonumber
\end{align}
They express the exact identity
\[
 h_R+h_I\sqrt{B_*}=p^8\mathcal H_8((W+\sqrt{B_*})/p).
\]
Taking its conjugate product gives \eqref{eq:resultant}, since the quadratic's
leading coefficient is $p^2$. No division or square root is used in the
actual evaluator. The complete 117-gate source is
\pathtext{data/merger12_product.json}; the equations above are its mathematical
specification, not a hidden extra-variable system.

\section{Why the elimination is sound}\label{sec:proof}

\subsection{The monic-root merger}
\begin{lemma}\label{lem:monic}
Let $a,c\in\pos$, $\Delta=(a+2)^2-1$, $H=4a+3$, and
$B_*=1+\Delta c^2$. Suppose a monic polynomial in $\Z[v]$ has a common root
with $(W-Hcv)^2-B_*$, where $W\in\Z$. Then $B_*$ is a square and the common
root is an integer.
\end{lemma}
\begin{proof}
The common root is $(W\pm\sqrt{B_*})/(Hc)$ and is an algebraic integer by
monicity. If $B_*$ is nonsquare, the difference of this root and its conjugate
is a nonzero algebraic integer in the quadratic field. Its norm is
\[
 -\frac{4B_*}{H^2c^2}.
\]
But
\[
 H^2-4\Delta=12a^2+8a-3\geq17,
\]
so $0<4B_*<H^2c^2$. The norm lies strictly between $-1$ and $0$, contradicting
the integrality of the norm of an algebraic integer. Thus $B_*$ is square.
The root is now rational, and every rational algebraic integer is an integer.
\end{proof}

Applied to \eqref{eq:h8}, this lemma gives an integer $v$. Consequently
\begin{equation}\label{eq:restoregamma}
 \gamma=u-cv\in\Z,
 \qquad (X+ac+H\gamma)^2-\Delta c^2=1.
\end{equation}
It is essential that $\gamma$ may initially be negative. The inherited
signed-domain theorem handles that sign; we do not assume the positive main
root merely because the original compiler used one.

\subsection{A rational Pell ordinate saturates}
\begin{lemma}\label{lem:saturation}
If $S\geq2$ and $V$ are integers and $y\in\mathbb Q$ satisfies
\[
 S^2V^2-(S^2-1)y^2=1,
\]
then $y\in\Z$.
\end{lemma}
\begin{proof}
Write $S^2-1=b^2D_0$ with $D_0>1$ squarefree. Put
$x=S|V|$ and $z=b|y|$. Then $x^2-D_0z^2=1$. In a reduced rational
expression for $z$, its denominator squared must divide $D_0$; hence $z$ is
integral. Neither $V$ nor $y$ is zero.

Let $F+G\sqrt{D_0}$ be the fundamental positive norm-one unit with integer
coordinates. Ordinary Pell classification gives positive indices $e_0,m$ with
\[
 S=\chi_F(e_0),\quad b=G\psi_F(e_0),\quad
 x=\chi_F(m),\quad z=G\psi_F(m).
\]
Here the sequences are normalized by
$(F+\sqrt{F^2-1})^n=\chi_F(n)+\psi_F(n)\sqrt{F^2-1}$.
Because $S\mid x$, we have $\chi_F(e_0)\mid\chi_F(m)$.

The identities
\[
 \chi_F(j+2e_0)\equiv-\chi_F(j),\qquad
 \chi_F(2e_0-j)\equiv-\chi_F(j)
 \pmod{\chi_F(e_0)}
\]
reduce $m$ to an index $0\leq j\leq e_0$. Strict growth excludes every
$j<e_0$, since $1\leq\chi_F(j)<\chi_F(e_0)$. Therefore $m$ is an odd
multiple of $e_0$. The multiplication identity implies
$\psi_F(e_0)\mid\psi_F(m)$, and
$|y|=\psi_F(m)/\psi_F(e_0)$ is integral.
\end{proof}

This is stronger than the generally false inference that a rational number
whose square is an arbitrary rational number must be integral. The special
square coefficient $S^2$ is what makes the lemma work.

\subsection{Restoration of the other three roots}
For a rational $t$, define
\[
 C_t=dt^2+N-e,\qquad
 U_t=C_t^2+4(a_0-dt^2N),\qquad
 \mathcal E(t)=U_t^2-16C_t^2a_0.
\]
Direct expansion gives the polynomial identity
\begin{equation}\label{eq:monicidentity}
 \mathcal H_8(dt)=d^4\mathcal E(t).
\end{equation}

We first justify the denominator bounds without assuming any deleted root.
When $N_k,N_t$ are equal units and the twelve supplied variables are positive,
weak transport and \eqref{eq:bootstrap} hold. Since $h\geq1$, we have $k>E$.
Put $A_0=a+2$. Then $A_0<2E$, $c<2kY$, and
\[
 L_0>A_0c/4,\qquad L_0>4A_0^2,\qquad
 D<A_0^2I^2,\qquad b_0<2kA_0^2I^2.
\]
For $M_0=kI^2+T^2$,
\[
 M_0^2\geq4kI^2T^2,\qquad M_0^2\geq k^2I^4.
\]
Together with $A>L_0^2$, these bounds give
$AM_0^2>4Dv_0^2$ and $AM_0^2>4b_0^2$. For the first inequality it is enough
that $k>16$, which follows already from $k>E$. Splitting into two halves gives
\[
 AM_0^2>2(Dv_0^2+b_0^2).
\]
Since $d=Q-1\geq Q/2$ and $M=2M_0$, we obtain the strict bound
\begin{equation}\label{eq:domination}
 N=dM^2A>4(e+1)>0.
\end{equation}

Suppose now $\mathcal E(t)=0$ for rational $t$. We have $C_t>0$. Also $t\ne0$,
because
\[
 \mathcal E(0)=((N-e)^2-4a_0)^2>0,
 \qquad 2\sqrt{a_0}\leq P+Z=e+1<N-e.
\]
Put $B_0=2Qv_0b_0>0$. Then $4a_0=DB_0^2$, and
\[
 f=-\frac{U_t}{2C_tB_0}
\]
is rational with $f^2=D$, hence an integer. The definition of $U_t$ now yields
\[
 (C_t+B_0f)^2=4d^2t^2M^2A.
\]
Thus
\[
 \tau=\frac{C_t+B_0f}{2dtM},\qquad y=t-M\tau
\]
give an integer $\tau$ with $\tau^2=A$ and initially rational $y$. Expanding
the defining relations gives
\[
 Q(v_0f-b_0)^2-(Q-1)y^2=1.
\]
Here $v_0f-b_0=c(Tf-1)-Rf^2$ is integral, and
$Q=(\Delta i c^2)^2$. Lemma~\ref{lem:saturation} therefore gives $y\in\Z$.
It also follows that $t=M\tau+y$ is integral. The strong equation follows
from $f^2=D$. The restored first and auxiliary roots are nonzero.

\subsection{System and single-polynomial soundness}
At a three-equation solution, $p>0$ and $\mathcal H_8$ is monic, so the
vanishing resultant has a common root. Lemma~\ref{lem:monic} and
\eqref{eq:restoregamma} restore the main norm. Applying the previous argument
to $t=v/d$ restores $f,\tau,y$. The signed ordinary interface, using
\eqref{eq:ordinaryconvert}, now gives a blank halting computation.
One may change both signs $(f,T)$ together and replace $\tau,y$ by their
absolute values; no sign change is assumed to preserve an expression that
depends on an unsquared root.

For the single equation, $N_k,N_t,1+\mathcal R_{12}$ are integer units.
The formula \eqref{eq:normcircuit} has even $h_I$, so
$\mathcal R_{12}\equiv h_R^2\pmod4$. Thus $1+\mathcal R_{12}\ne-1$,
and $\mathcal R_{12}=0$. The two outer factors have the same sign. The
preceding bounds explicitly allowed both signs. After restoration, the
inherited signed Pell step-down forces $N_k=N_t=1$ and $C>0$.
In particular, the exceptional boundary $C=0$ has not been excluded
prematurely.

\subsection{Completeness}
Start with any canonical positive normalized compiler witness for a finite
blank halting run. Set
\[
 t=M\tau+y,\qquad v=dt,\qquad u=\gamma+cv.
\]
All twelve supplied coordinates are positive. The original norm identities
give \eqref{eq:monicidentity} with $\mathcal H_8(v)=0$, and the main norm
gives $(W-pv)^2=B_*$. Hence $\mathcal R_{12}=0$, and all three product
factors equal $1$. This proves the converse implication.

\section{Exact size and the thirteen-variable alternative}\label{sec:degree}

\subsection{The degree is 1529, not a gate-propagation guess}
Give every supplied coordinate degree one and the four coefficients degree
zero. The relevant exact degrees are
\begin{center}\small
\begin{tabular}{lrrrrrrrrr}\toprule
Expression &$c$&$a$&$\Delta$&$R$&$A$&$I$&$D$&$Q,d$&$P$\\
Degree&5&6&12&4&22&11&34&46&92\\\midrule
Expression &$Z,e$&$a_0$&$M$&$N$&$B_*$&$W,p$&$v_R$&$h_R$&$h_I$\\
Degree&122&214&23&114&22&11&380&760&581\\\bottomrule
\end{tabular}
\end{center}
The unique highest term of $v_R$ is $d^2p^4e^2$; the highest term of $h_R$
is $d^4p^8e^4$. Therefore $\mathcal R_{12}$ has the nonzero highest form
$d_{\rm top}^8p_{\rm top}^{16}e_{\rm top}^8$, of degree 1520; the subtracted
norm term has degree at most 1184. If $q_0=\beta J$ and $k_0=\eta+\zeta$,
this highest form is explicitly
\[
 4^{16} i^{64}k_0^{144}w^{112}s^{256}q_0^{928}(q_0-F)^{16}.
\]
It is nonzero for every authentic $\beta>0$. The leading forms of $N_k$ and
$N_t$ have degrees 7 and 2 and are also nonzero. Their product proves
$\deg\PS_{12}=1529$. The emitted circuit costs $73$ multiplications and
$44$ additions/subtractions, totaling $117$. Deleting the product finalizer
leaves $112$ operations for the three equations when the final norm terms
are compared directly.

\subsection{Retaining the strong root lowers degree substantially}
The twelve variables in \eqref{eq:variables}, together with $f$, give
thirteen variables.
Retain the common outer definitions, and now put
\begin{align*}
 I&=ic^2,&K_s&=\Delta(f^2-1),&d_s&=K_s-1,\\
 V_f&=c(Tf-1)-Rf^2,&e_s&=K_sV_f^2-1.
\end{align*}
The monic quartic
\[
 \mathcal H_4(v)=(v^2+d_s^2A-d_se_s)^2-4d_s^2Av^2
\]
replaces \eqref{eq:h8}. Define
\[
 \mathcal R_{13}=\Res_v(\mathcal H_4(v),(W-pv)^2-B_*).
\]
The system is $N_k=N_t=1$, $I^2=K_s$, $\mathcal R_{13}=0$, and the single
polynomial is
\begin{equation}\label{eq:thirteen}
 \boxed{\PS_{13}=N_kN_t(1+I^2-K_s)(1+\mathcal R_{13})-1.}
\end{equation}

For an explicit norm evaluator, use $r_0=W^2+B_*$ and $r_1=2W$, and put
\begin{align*}
 s_R&=r_0+p^2d_s(d_sA-e_s),\\
 h_R&=s_R^2+B_*r_1^2-4d_s^2Ap^2r_0,\\
 h_I&=r_1(2s_R-4d_s^2Ap^2),\\
 \mathcal R_{13}&=h_R^2-B_*h_I^2.
\end{align*}
The same monic lemma gives integral $v$ and $\gamma=u-cv$. Put $t=v/d_s$.
The quartic identity restores
$\tau=(d_st^2+d_sA-e_s)/(2d_st)$ and $y=t-\tau$.
Here $t=0$ is impossible: that branch would give $K_s\mid A-1$, whereas
the weak outer bounds $k>E$, $Y\geq2$ imply
\[
 0<A-1=k^2EY(EY+1)<k^4Y^4<c^4\leq I^2=K_s.
\]
The rational-square argument restores $\tau$, and Lemma~\ref{lem:saturation}
restores $y$, since $K_s=I^2$. The signed ordinary interface proves soundness.
Completeness uses $u=\gamma+cd_s(\tau+y)$.

The strong guard cannot be $-1$: $K_s$ is $0$ or $1$ modulo $4$, as is
$I^2$, so $K_s-I^2$ cannot be $2$ modulo $4$. The resultant guard has the
same even-imaginary-part argument as before. Thus the product has exactly
the required existence condition.

The highest form of $h_R$ is $d_s^4A^2p^4$, of degree 144. Hence
$\deg\mathcal R_{13}=288$, and
\[
 \deg\PS_{13}=288+22+7+2=319.
\]
The system costs 83 operations and the product 91. This is a useful
intermediate choice when degree matters more than eliminating the last variable.

\section{Twenty-six unrestricted integer unknowns}\label{sec:integers}

The twelve-variable formula has a larger sound signed domain. Keep exactly
\[
 J,F,\alpha,z,s,\eta,\zeta
\]
positive. Let $h,T,u,r,v_w$ be unrestricted integers and substitute
\begin{equation}\label{eq:signedsub}
 i=2r+1,\qquad w=v_w^2.
\end{equation}
Thus the underlying formula has seven positive and five signed coordinates.
Replace each positive coordinate $x$ by its own disjoint copy of
\begin{equation}\label{eq:pos}
 \operatorname{Pos}(a_x,b_x,c_x)=1+a_x^2+b_x^2+c_x(c_x+1).
\end{equation}
This is always a positive integer and represents every positive integer.
Indeed, the three-square theorem applied to $4x-3$ gives one odd square and
two even squares; dividing the corresponding identity by four gives
\eqref{eq:pos}. The resulting polynomial $\PS_{26,\Z}$ has
\[
 \boxed{7\cdot3+5=26}
\]
unrestricted integer unknowns. No positivity condition remains outside the equation.

\subsection{Why the enlarged domain is sound}
The product unit argument still forces $\mathcal R_{12}=0$. Even if $w=0$,
we have $a=Y>0$, $c>0$, and $i\ne0$, so the monic lemma applies first and
restores an integral $\gamma=u-cv$ and the main norm. If $w=0$, then
$X=0$ and reduction of this norm modulo $H$ would give $0\equiv1\pmod H$.
Hence $w>0$ and the usual outer bounds follow.

The index equation with a possible unit sign forces $h\geq0$, since
$h\leq-1$ would imply $k\leq R+1-E<0$. Thus $k\geq R-1$ and $R<2k$.
The old estimate $k>E$, valid in the positive proof, may not be reused here.
The following replacement supplies the necessary domination.

The restored main norm classifies $c=\psi_{A_0}(j)$, where $A_0=a+2$.
Since $c>1$, $j\geq2$. The case $j=2$ gives $c=2a+4$ and main root
$\pm(2a^2+8a+7)$. Its congruence modulo $H$ forces respectively
$X\equiv4$ or $X\equiv3a+2\pmod H$, impossible because $16\leq X<a<H$.
Therefore $j\geq3$, so
\[
 c>4a^2,\qquad k>2X^2Y,
 \qquad L_0>12A_0^2,\qquad L_0>A_0c/4.
\]
These imply $A>144A_0^4$ and $A>A_0^2c^2/16$. With
$M_0=kI^2+T^2$, $D<A_0^2I^2$, and $b_0<3kA_0^2I^2$, one obtains
\[
 AM_0^2>12Dv_0^2,\qquad AM_0^2>16b_0^2,
 \qquad AM_0^2>4(Dv_0^2+b_0^2).
\]
This is more than enough for the doubled weight in the actual formula.
In particular $C_t>0$ and $t=0$ is impossible even when $t,T,u$ are signed.

There is one additional denominator case: $T=0$. Then $v_0=0$ and the
eliminant reconstructs $\tau=C_t/(2dtM)$ with $\tau^2=A$, hence integer
$\tau$, and a rational $y$ with
\[
 Qb_0^2-(Q-1)y^2=1.
\]
Lemma~\ref{lem:saturation} makes $y$ integral. But $1<b_0<Q$: indeed
\[
 Q-b_0=(\Delta-R)D-\Delta-c
 \geq1+\Delta(i^2c^4-1)-c>0.
\]
There is no integer solution to this norm in that interval. Its rearrangement
$(Q-1)(y^2-b_0^2)=b_0^2-1>0$ implies $|y|\geq b_0+1$ and therefore
\[
 b_0^2-1\geq(Q-1)(2b_0+1)\geq b_0(2b_0+1)>b_0^2-1,
\]
a contradiction. Thus $T\ne0$. All denominators in root restoration are now
nonzero, and the same algebra gives integral $f,\tau,y,\gamma$.
Package 14's signed normalized interface applies with precisely its seven
positive coordinates, odd $i$, and square $w$.

\subsection{Completeness and size}
Canonical compiler tableaux can have width and height powers of five. Since
$d_{\rm base}=5^{16}$ and their cyclic length are odd, $q=2^{d_{\rm base}N}$
has odd exponent. The recovered rank $R$ is odd, so
$w=2^{R-d_{\rm base}N}$ is a positive square. Choosing the strong index
$m=cR$ yields an odd $i$, as in the inherited signed-domain proof. This
preserves existence; it does not claim every old witness already has these
properties. The positive merger witness $u=\gamma+cdt$ is allowed in the
larger signed domain. Finally use the surjections \eqref{eq:pos}.

The seven positivity replacements cost seven operations each; odd $i$ costs
two and square $w$ one. Thus the product costs $117+49+3=169$ operations,
and the system costs 164. The exact degree of the product is \IntegerDegree;
the system's resultant has degree 2976. To see this, assign weight two to
the seven positive ports and $w$, weight one to $h,i,T,u$, and weight zero
to the fixed coefficients. The top degrees are
\[
 c=10,\ a=12,\ \Delta=24,\ d=90,\ e=238,\ p=W=22.
\]
The norm's unique leading form has degree
$8(90)+16(22)+8(238)=2976$. Adding the outer degrees $13$ and $4$ gives
2993. Substitution by the disjoint positive quadratic forms does not cancel
this leading form: setting two coordinates in each triple to zero reduces
it to an injective squaring substitution.

\section{A 59-operation system}\label{sec:system59}

There is a small independent improvement for the arithmetic-count objective.
Retain the sixteen positive variables
\[
 J,F,\alpha,z,f,h,i,o,j,s,w,\tau,\eta,\zeta,y,\gamma.
\]
Use the common definitions and put
\[
 I=ic^2,\qquad K_s=\Delta(f^2-1),\qquad V_f=of-c,\qquad U=jc-R.
\]
The seven polynomial equations are
\begin{align}
 \tau^2-L_0(L_0+k)&=1,\nonumber\\
 (X+ac+H\gamma)^2-\Delta c^2&=1,\nonumber\\
 K_s(V_f^2-y^2)+y^2&=1,\nonumber\\
 k-hE-R&=1,\label{eq:s59}\\
 (K_0+w)C&=F+z(q-1),\nonumber\\
 I^2&=K_s,\nonumber\\
 U&=V_f.\nonumber
\end{align}
The coordinates $o$ and $j$ retain the two auxiliary divisibility conditions;
there is no supplied $T$ coordinate in this system. This is the retained
ordinary-strong interface proved in package 09. Its circuit is supplied in
\pathtext{data/retained16_system59.json}.

The actual improvement is the coordinate change $z_{\rm old}=z+1$.
In the old transport equation
\[
 (K_0+w)C+(q-F)-z_{\rm old}(q-1)=1,
\]
this gives $(K_0+w)C=F+z(q-1)$ and saves one binary operation. At any old
positive zero, the compiler gives $w>q>F$ and $C>0$, so $z_{\rm old}=1$
would force $(K_0+w)C=F$, a contradiction. Hence $z=z_{\rm old}-1>0$.
Conversely every new positive solution gives an old one by adding one.
The map is a bijection of positive zero sets.

The emitted system has 33 multiplications and 26 additions/subtractions,
for 59 operations. Its seven residual degrees are $22,18,24,7,2,22,6$.
A sum of their squares gives a single degree-48 equation using 79 operations.
This system is preferable when evaluation cost matters more than the number
of supplied integers. It does not improve the 67-operation single polynomial.

\section{A smaller fixed coefficient recipe}\label{sec:coefficients}

\subsection{Literal machine and cell alphabet}
The machine is the accepted dense K215 table, with working states $0$ through
$214$, blank symbol $0$, start state $0$, and halt state $215$. The archive
contains its literal table, with SHA-256
\begin{center}\footnotesize\ttfamily
fdbcef0825886035c3091c8537a69535bc96f9e265f0ae45863543d40ae70090.
\end{center}
The actual first action is $(0,0)\mapsto(1,+1,93)$. All original actions,
including the movement into halt, are retained. Only the padding transitions
$(215,b)\mapsto(b,0,215)$ are added after halting.

There are 439 cell symbols. Headless bits have codes $0,1$; a bit $b$ carrying
state $q$ has code $2(q+1)+b$, for $0\leq q\leq215$. Initial phases have
codes $L=434,Q_0=435,R_0=436$, carrying blank payloads and a head only at
$Q_0$. A vertical wall is $V=437$ and a horizontal time boundary is $H_0=438$.
The notation $Q_0,R_0,H_0$ in this section is unrelated to the arithmetic
abbreviations $Q,R,H$ above.

A tile is a four-cell tuple $(l,c,r,n)$, with $n$ the next-time center.
Begin with the exact T4 predicate of package 11: at most one head in
$(l,c,r)$, propagation and nonadjacency of walls, horizontal propagation of
$H_0$ inside a strip, initialization after $H_0$, phase edges
$LL,LQ_0,Q_0R_0,R_0R_0$, no head escape, exact deterministic update, and
halting before the next $H_0$.

Strengthen that predicate by the following restrictions:
\begin{enumerate}
\item Both neighboring pairs $(l,c)$ and $(c,r)$ obey the horizontal
wall/phase restrictions. For a pair not touching $V$, both entries have the
same $H_0$ status and the same phase status; phase pairs are in the four
allowed edges. Across a wall, the adjoining phase must be $L$ on its right
and $R_0$ on its left; two neighboring walls are forbidden.
\item If a horizontal triple contains $V$, all its symbols belong to
$\{0,L,R_0,V,H_0\}$. Thus nearby ordinary cells are blank and headless.
\item If $n=H_0$, none of $l,c,r$ is an initial phase and every head in the
triple is in halt state.
\item If $c=H_0$ and $l=V$, require $n=L$; if $r=V$, require $n=R_0$.
\end{enumerate}
There are exactly 5,251 tuples satisfying the strengthened rule. Add the
old locally valid but globally impossible tile
$E_0=(H_0,0,0,H_0)$, and order the alphabet as
\[
 S_0=(L,Q_0,R_0,1),\quad E_0,\quad
 \text{all remaining strengthened tuples in lexicographic order}.
\]
The resulting list has 5,252 entries and is supplied in
\pathtext{data/strict_windows.json}. It is a finite list computed from the
430 literal actions; generating it does not search for a K215 halt.

\subsection{Why the stricter language preserves the target computation}
Every tile in the new list is an old T4 tile. Therefore any periodic new
tableau gives an old one, and the established marked-strip argument proves
soundness. The forbidden tile $E_0$ cannot occur globally: its left neighbor
is $H_0$ with a nonwall right neighbor $0$, contradicting the horizontal
boundary rule at that left cell. It is retained only as the selector that
the arithmetic mask excludes.

For completeness, pad a finite genuine halting run with at least two
unvisited blank cells next to each wall. Such cells remain blank throughout.
Every ordinary triple touching a wall is then blank and headless; initial
rows use $L$ and $R_0$ at the corresponding margins, and boundary rows use
$H_0$. All other new restrictions hold on genuine rows. In the helical
periodic presentation, a wall separates neighboring strips whose time
origins differ by one. The allowed wall alphabet includes every combination
of their possible blank, phase, and boundary edge symbols, so these seams
also satisfy the restrictions. Independent width and height padding remain
available, including arbitrarily large powers of five.

The original three overlaps are unchanged:
\[
 w_i(r)=w_{i-1}(c),\qquad w_i(c)=w_{i-1}(l),\qquad
 w_i(n)=w_{i-h}(c).
\]
They reconstruct the plane from the center word even when residues collide.
No condition $h\mid N$ or $\gcd(h,N)=1$ has been added.

\subsection{Exact four integers}
Set
\[
 a_c=439,\ k_c=5252,\ \ell=13,\ A_c=2^{13},\
 \rho=2^{78125},\quad B_{\mathrm{base}}=\rho^{1953125}=2^{5^{16}}.
\]
For $0\leq j\leq438$, the four copy positions are
\[
 e_l(j)=8764+j,\quad e_c(j)=8325+j,\quad
 e_r(j)=7886+j,\quad e_n(j)=7008+j.
\]
Let the native support $\mathcal S$ consist of selectors $0,\ldots,5251$,
these four copy intervals, dummy positions $10520,10521$, and anchors
\[
 M_0=11839,\quad3M_0=35517,\quad5M_0=59195,\quad11M_0=130229.
\]
Thus $|\mathcal S|=7014$. Set
\[
 \mu=A_c-2+\sum_{j=1}^{2835}A_c^j.
\]
For tile $w_i=(l_i,c_i,r_i,n_i)$ define
\begin{align*}
 v_i={}&1+A_c^{1+l_i}+A_c^{440+c_i}+A_c^{879+r_i}
          +A_c^{1318+n_i}+\sum_{t=1}^{4}A_c^{1756+t},\\
 v_{e_l(j)}={}&A_c^{1+j},&v_{e_c(j)}={}&A_c^{440+j},\\
 v_{e_r(j)}={}&A_c^{879+j},&v_{e_n(j)}={}&A_c^{1318+j}.
\end{align*}
At anchors $M_0,3M_0,5M_0,11M_0$, respectively, use
$A_c^{1757},A_c^{1758},A_c^{1759},A_c^{1760}$. Both dummy coefficients are
zero. Put
\begin{align}
 D_C={}&\rho^{1317}+\rho^{226698}+\rho^{47356}+\rho^{94712}
       +\sum_{e\in\mathcal S}v_e(\rho^{487157-e}+\rho^{747619-e}),\nonumber\\
 D_R={}&\rho^{226259},\nonumber\\
 M_F^0={}&\mu(\rho^{487157}+\rho^{747619})
       +\sum_{j=0}^{438}\rho^{8325+j}\label{eq:coefficientrecipe}\\
       &+\sum_{j=0}^{438}(\rho^{226259+8764+j}+\rho^{226259+8325+j})
       +\rho^{59195}+\rho^{130229},\nonumber\\
 M_C={}&B_{\mathrm{base}}-1-
          \sum_{\substack{e\in\mathcal S\\e\ne1}}\rho^e-2\rho^{10520},
 \qquad M_F=M_F^0+4.\nonumber
\end{align}
The four coefficients are now exactly
\[
 \boxed{\beta=B_{\mathrm{base}}-1,\quad K_0=D_C+B_{\mathrm{base}}D_R,
 \quad M_C,\quad M_F^*=M_F+B_{\mathrm{base}}-1.}
\]
All exponentiations in this recipe have fixed integer exponents. These
constants are independent of the unknowns and of whether K215 halts.

\subsection{Compiler geometry and population proof}
The new anchor set is $\{M_0,3M_0,5M_0,11M_0\}$. The tested temporal
anchor targets are $5M_0,11M_0$, reached by bare shifts $4M_0,8M_0$
from the first two anchors. The only support pair with difference $6M_0$
is $(5M_0,11M_0)$: nonanchor support is below $M_0$, and the other anchor
differences are different. The circumference is larger than twice the
support diameter, so there is no wrapped alternative. This replaces the
older anchor set $\{M_0,3M_0,9M_0,27M_0\}$.

The remaining scalar layout is
\[
 E_{\max}=130229,\quad H_c=226259,\quad
 T_1=487157,\quad T_2=747619,\quad L_c=1953125.
\]
All copy/anchor targets lie below $T_1-E_{\max}$. At each target, the
four bare shifts in $D_C$ hit exactly the intended source positions, and
$D_R=\rho^{H_c}$ tests exactly the intended spatial copies. There are
1,319 such targets. The supplied finite set checks exhaust all of them.
The clause bands have separation $T_2-T_1=260462>2E_{\max}$, while the
opposite circular separation is also larger than $E_{\max}$. Thus every
rotation leaves at least one clean clause band.

All bits cleared from $M_C$ are distinct and some highest bits remain;
all fields of $M_F$ lie strictly below bit $5^{16}$ and leave unused high
bits. Thus $0<M_C,M_F<B_{\mathrm{base}}-1$. The selector at position zero
clears bit zero of $M_C$, while bit one remains. The only contribution to
$M_F$ below the first radix position is the marker $4$. Therefore
$M_C\equiv2\pmod4$ and $M_F\equiv4\pmod8$.

The mask population is exact:
\[
 \pc(\mu)=12+2835=2847,
 \quad\pc(M_F)=2(2847)+3(439)+2+1=7014,
 \quad\pc(M_C)=5^{16}-7014.
\]
The last $+1$ in $\pc(M_F)$ comes from the native low-bit marker $4$.
There are two low dummy positions, both with zero clause coefficient. Only
the first receives the extra allowed upper bit through $-2\rho^{10520}$
in $M_C$. Hence there is still only one exceptional upper spill digit in
the all-rotation proof. The two anchor tests exclude every nonzero binary
rotation within a radix digit; the unique anchor difference then excludes
every nonzero radix rotation within a cell.

The carry bounds are also checked with manageable integers:
$\mu$ has bit length 36,856, the sum of clause coefficients has bit length
22,893, and the required mass threshold has bit length 22,909. The chosen
radix exponent 78,125 is larger than all of these bounds, and the native
compiler's carry-free estimates apply.

Finally, $b=78125$ is odd, so $\rho\equiv2\pmod5$.
The band separation is $2$ modulo $4$, so paired clause contributions in
$D_C$ cancel modulo $5$. The two new bare anchor shifts are multiples of
$4$, and $439\equiv3\pmod4$. Therefore
\[
 2D_C-D_R\equiv3\pmod5.
\]
This is the unit required to tune the actual packed rank using the first
dummy bit. Width, height, and $d_{\rm base}$ remain powers of five.
The original dummy-control completeness and canonical positivity margin
therefore apply unchanged. Together with the local-language proof, this
establishes the complete four-port interface for the new tuple.

\section{Metatheory, verification, and integration}\label{sec:verification}

\subsection{The logical endpoint}
The existing K215 endpoint is
\[
 \operatorname{Halt}(K215,\mathrm{blank})
 \longleftrightarrow
 \text{there is a finite ZFC derivation of falsity}.
\]
Negating this and composing with the proved Diophantine equivalence gives
Theorem~\ref{thm:main}. No consistency premise is used to prove the
biconditional. If ZFC is inconsistent, both a proof-search halt and a finite
integer witness exist. If it is consistent, the corresponding equation has
no solutions in the stated domain.

The distinction between this statement and a proof of consistency is
substantive. Under the usual second-incompleteness hypotheses, consistency
prevents ZFC from proving the no-solution statement. Mere consistency alone
does not justify a claim of two-sided unprovability; an appropriate soundness
assumption is needed for the stronger claim. No two-sided independence claim
is required here.

For ordinary nonnegative natural-number unknowns, replace each positive
coordinate of $\PS_{12}$ by $x+1$. This preserves the number of unknowns
and the degree. It adds twelve operations to the literal evaluator. Do not
interpret the original positive polynomial directly over all of $\Z$;
Section~\ref{sec:integers} supplies the separately proved unrestricted form.

\subsection{What the checks establish}
The archive contains division-free JSON circuits, standard-library Python
emitters, and recorded checks. The checks establish the live operation counts,
algebraic identities, degree bounds and leading-form witnesses, agreement
with an independent quadratic-pair evaluation, the literal machine/table
binding, exact finite tile enumeration, field positions, carry bounds,
anchor synchronization, and padded control-tableau seams.

The checks do not establish $\Con(\mathrm{ZFC})$, do not produce an accepting
K215 run, and do not replace the unbounded proofs in this report. The
enormous coefficients are never expanded. No fresh end-to-end Lean or Rocq
build was performed. Independent reviews of the new monic elimination and
strict compiler are included in \pathtext{proofs/}.

\subsection{Reproduction and file layout}
Run the included reproduction driver with Python 3.10 or later. It uses
only the standard library for the required arithmetic and finite-data
checks. Compile this report with pdfLaTeX twice, or with the provided
Makefile. The top-level README gives the exact commands and integration
destination. The archive is self-contained for its new checks and includes
the inherited mathematical interfaces used in the proofs.

The recommended integration is a new numbered synthesis directory under
the existing \pathtext{con-zfc-diophantine-forms} report collection. Keep the
earlier packages as comparison points for objectives they still optimize.
In particular, retaining the 67-operation polynomial is appropriate even
though its coefficient tuple can be rebound to the new strict compiler.

\subsection{Remaining work}
A useful next formalization target is Lemma~\ref{lem:saturation}, followed
by Lemma~\ref{lem:monic} and the monic identities. These are relatively small
arithmetic results and isolate the genuinely new proof obligations. The
larger task remains formalization of the native-window/Pell bridge.
Further variable elimination would need a new integrality argument for
another quotient; simply clearing its denominator or taking another
resultant would not establish the required equivalence.

\clearpage
\begin{thebibliography}{9}
\bibitem{Repo} V. Reshetnikov et al., \emph{ProveIt}, snapshot \pin,
22 October 2026. Repository paths are given in the package's provenance
manifest and in \pathtext{reference/comparison_at_pin.md}.\newline
\url{https://github.com/VladimirReshetnikov/ProveIt}.
\bibitem{Wave2} ProveIt, \emph{Con(ZFC) in Diophantine form}, wave-2 packages
09--15, at the pinned snapshot. Principal dependencies: 09 weighted roots;
11 T4 compiler; 12 ordinary signed interface; 14 signed-domain reduction.
\bibitem{K215} ProveIt, \emph{Formal correctness of the exact 215-state machine},
and \emph{candidate-k215/proof.md}. The retained endpoint names are
\pathtext{KEndpoint215.blank_halts_iff_zfc_refutation} and
\pathtext{tm215_halts_iff_ZFC_inconsistent}.
\bibitem{Jones} J. P. Jones, \emph{Universal Diophantine equation},
Journal of Symbolic Logic \textbf{47} (1982), 549--571.
\url{https://doi.org/10.2307/2273588}.
\bibitem{MR} Yu. Matiyasevich and J. Robinson,
\emph{Reduction of an arbitrary Diophantine equation to one in 13 unknowns},
Acta Arithmetica \textbf{27} (1975), 521--553.
\url{https://doi.org/10.4064/aa-27-1-521-553}.
\bibitem{Bayer} J. Bayer, M. David, M. Hassler, Yu. Matiyasevich, and
D. Schleicher, \emph{Diophantine Equations over $\Z$: Universal Bounds and
Parallel Formalization}, arXiv:2506.20909v2 (2025).
\url{https://arxiv.org/abs/2506.20909}.
\end{thebibliography}
\end{document}
